\clearpage\pagestyle{otMain}\setlength{\otBlocEspacement}{2.5mm}\fontsize{11}{15}\selectfont\corrigesChapitre{17}{Symétrie axiale}
\begin{blocLivre}{cor-17-1}\begin{corrige}{1}{Combien d’axes ?}\demarche{Imaginer un pliage qui superpose les deux moitiés.}
Rectangle : $\res2$. Carré : $\res4$. Triangle isocèle : $\res1$. Losange : $\res2$, ses diagonales. Triangle quelconque : $\res0$. Les axes sont tracés en vert.
\label{fig-cor-17-1}\figureGeo[corrige]{sym-1}
\end{corrige}\end{blocLivre}
\begin{blocLivre}{cor-17-2}\begin{corrige}{2}{Axe vertical}\demarche{Compter perpendiculairement à l’axe.}
On compte les carreaux horizontalement jusqu’à l’axe et on les reporte de l’autre côté. L’image et le point $A'$ figurent en vert.
\label{fig-cor-17-2}\figureGeo[corrige]{sym-2}
\end{corrige}\end{blocLivre}
\begin{blocLivre}{cor-17-3}\begin{corrige}{3}{Axe horizontal}\demarche{Compter perpendiculairement à l’axe.}
On compte les carreaux verticalement : chaque point est reporté à la même distance, de l’autre côté de l’axe. L’image et $B'$ figurent en vert ; le drapeau se retrouve tête en bas.
\label{fig-cor-17-3}\figureGeo[corrige]{sym-3}
\end{corrige}\end{blocLivre}
\begin{blocLivre}{cor-17-4}\begin{corrige}{4}{Axe en diagonale}\demarche{Suivre les diagonales des carreaux.}
On avance en diagonale à travers les carreaux, perpendiculairement à l’axe. Les écarts horizontaux et verticaux s’échangent ; l’image et $C'$ figurent en vert.
\label{fig-cor-17-4}\figureGeo[corrige]{sym-4}
\end{corrige}\end{blocLivre}
\begin{blocLivre}{cor-17-5}\begin{corrige}{5}{Papier pointé}\demarche{Les points jouent le rôle des carreaux.}
On compte les intervalles entre les points, perpendiculairement à (d), et on les reporte. L’image est en vert ; les longueurs sont conservées.
\label{fig-cor-17-5}\figureGeo[corrige]{sym-5}
\end{corrige}\end{blocLivre}
\begin{blocLivre}{cor-17-6}\begin{corrige}{6}{Symétrique d’un point}\demarche{Perpendiculaire, puis report de longueur.}
On trace la perpendiculaire à (d) passant par A ; elle coupe (d) en H. On reporte la longueur AH au compas de l’autre côté de (d) : c’est $A'$, en vert.
\label{fig-cor-17-6}\figureGeo[corrige]{sym-6}
\end{corrige}\end{blocLivre}
\begin{blocLivre}{cor-17-7}\begin{corrige}{7}{Image d’un segment}\demarche{Construire les images des extrémités.}
On construit $A'$ et $B'$, puis on les relie (en vert). La symétrie conserve les longueurs : $\res{A'B'=4\ \mathrm{cm}}$.
\label{fig-cor-17-7}\figureGeo[corrige]{sym-7}
\end{corrige}\end{blocLivre}
\begin{blocLivre}{cor-17-8}\begin{corrige}{8}{Image d’un cercle}\demarche{Image du centre et même rayon.}
On construit le symétrique $O'$ du centre, puis le cercle de centre $O'$ et de même rayon, $\res{1{,}5\ \mathrm{cm}}$ (en vert).
\label{fig-cor-17-8}\figureGeo[corrige]{sym-8}
\end{corrige}\end{blocLivre}
\begin{blocLivre}{cor-17-9}\begin{corrige}{9}{Compléter la figure}\demarche{Compléter point par point.}
On construit le symétrique de chaque sommet par rapport à (d), puis on relie dans le même ordre. La moitié manquante est en vert.
\label{fig-cor-17-9}\figureGeo[corrige]{sym-9}
\end{corrige}\end{blocLivre}
\begin{blocLivre}{cor-17-10}\begin{corrige}{10}{Retrouver l’axe}\demarche{L’axe est la médiatrice d’un segment point-image.}
(d) est la médiatrice du segment $[AA']$. On la construit au compas ; elle est aussi la médiatrice de $[BB']$ et de $[CC']$, ce qui sert de contrôle. Elle figure en vert.
\label{fig-cor-17-10}\figureGeo[corrige]{sym-10}
\end{corrige}\end{blocLivre}
\begin{blocLivre}{cor-17-11}\begin{corrige}{11}{Ce que la symétrie conserve}\demarche{Utiliser les propriétés de conservation.}
La symétrie conserve longueurs, angles et aires : $\res{A'B'=5\ \mathrm{cm}}$, l’angle de sommet $B'$ mesure $\res{40^\circ}$ et l’aire vaut $\res{6\ \mathrm{cm}^2}$.
\label{fig-cor-17-11}\figureGeo[corrige]{sym-11}
\end{corrige}\end{blocLivre}
\begin{blocLivre}{cor-17-12}\begin{corrige}{12}{Les points de l’axe}\demarche{Un point de l’axe ne bouge pas.}
Le point I est sur (d) : il est son propre symétrique. Comme I est sur [MN], son image I est sur $[M'N']$ : le segment image passe par $\res I$. L’alignement est conservé.
\label{fig-cor-17-12}\figureGeo[corrige]{sym-12}
\end{corrige}\end{blocLivre}
\begin{blocLivre}{cor-17-13}\begin{corrige}{13}{Le billard}\demarche{Remplacer A par son symétrique.}
La symétrie conserve les longueurs : $AI=A'I$. Le trajet A-I-B a donc la même longueur que $A'$-I-B, qui est un segment. Le point d’impact I est en vert ; c’est le chemin le plus court qui touche la bande.
\label{fig-cor-17-13}\figureGeo[corrige]{sym-13}
\end{corrige}\end{blocLivre}
\begin{blocLivre}{cor-17-14}\begin{corrige}{14}{Lettres symétriques}\demarche{Chercher les pliages qui superposent.}
H : $\res2$ ; A : $\res1$, vertical ; E : $\res1$, horizontal ; N : $\res0$ ; X : $\res2$. Les axes sont en vert.
\label{fig-cor-17-14}\figureGeo[corrige]{sym-14}
\end{corrige}\end{blocLivre}
\begin{blocLivre}{cor-17-15}\begin{corrige}{15}{Un retournement}\demarche{Comparer l’orientation des deux figures.}
$\res{\text{Non}}$ : le drapeau pointait vers la droite, son image pointe vers la gauche. La symétrie retourne la figure, comme un miroir, tout en conservant ses longueurs et ses angles.
\label{fig-cor-17-15}\figureGeo[corrige]{sym-15}
\end{corrige}\end{blocLivre}
